\documentclass[10pt]{mypackage}

\usepackage[uprightgreeks,noDcommand]{kpfonts}
\renewcommand{\coloneq}{\coloneqq}
\renewcommand{\eqcolon}{\eqqcolon}
%\renewcommand{\mathbb}{\mathds}
%\usepackage{homework}
\usepackage{notes}

%\usepackage[ backend=bibtex, style = alphabetic, sorting=ynt ]{biblatex}
%\addbibresource{  }

\usepackage{parskip}

\fancyhf{}
\fancyhead[R]{Avinash Iyer}
\fancyhead[L]{Ergodic Theory: Class Notes}
\fancyfoot[C]{\thepage}

\setcounter{secnumdepth}{0}

\begin{document}
\RaggedRight
These are notes I'm taking for Ben Hayes's Ergodic Theory class for Fall 2026.

The main texts are Petersen's \textit{Ergodic Theory} and Walters's \textit{An Introduction to Ergodic Theory}, as well as some other 
\section{General Motivation and Common Examples}%
The basic structure of ergodic theory is a standard probability space $\left( X,\mu \right)$ and, in its most basic form, a measurable map $T\colon X\rightarrow X$ satisfying $T_{\ast}\mu = \mu$. Recall that the pushforward measure is defined by
\begin{align*}
  T_{\ast}\mu\left( A \right) &= \mu\left( T^{-1}(A) \right).
\end{align*}
We call $T$ a measure-preserving transformation.

Generalizing this, if $T$ is invertible, then we may take a group action $G\curvearrowright \left( X,\mu \right)$ by measure-preserving transformations. If $T$ is not invertible, then we may consider a semigroup action $N\curvearrowright \left( X,\mu \right)$.

The main items we are interested in are analytical considerations, such as
\begin{itemize}
  \item asymptotic structure
  \item connections to unitary representation theory
  \item isomorphism.
\end{itemize}
Ergodic theory appears in a number of fields, such as
\begin{itemize}
  \item unitary representation theory
  \item analytic and geometric group theory
  \item functional analysis
  \item probability theory
  \item geometry
  \item analytic number theory.
\end{itemize}
\subsection{Background}%
More aptly, this may be written as ``background.'' Most of these results are covered in a functional analysis class.
\begin{definition}
  A measurable space $\left( X,\Sigma \right)$ is called \textit{standard Borel} if there is a separable and completely metrizable space (also known as a Polish space) $K$ such that $\left( K, \mathcal{B} \right)\cong \left( X,\Sigma \right)$ are isomorphic as measurable spaces, where $ \mathcal{B} $ denotes the Borel $\sigma$-algebra on $K$.
\end{definition}
\begin{remark}
  In the special case of standard probability spaces, we may take $K$ to be compact. This emerges from the fact we are only concerned with Radon measures.
\end{remark}
\begin{example}
  The following are examples:
  \begin{itemize}
    \item finite or countably infinite sets;
    \item products of standard Borel spaces;
    \item $[0,1)\cong S^1$;
    \item $(0,1)\cong \R\cong [0,1)\times \Z$;
    \item $\R^{\N}$, where we use the complete distance metric
      \begin{align*}
        d\left( t,s \right) &= \sum_{j=1}^{\infty} 2^{-j} \frac{\left\vert t_j-s_j \right\vert}{1 + \left\vert t_j-s_j \right\vert}
      \end{align*}
      to provide the desired isomorphism;
    \item the set $\R\setminus \Q$ is a $G_{\delta}$ subset of the Polish space $\R$, so it is also a Polish space.
  \end{itemize}
\end{example}
\begin{theorem}
  If $\left( X,\Sigma \right)$ is standard Borel and $B$ is Borel, then $\left( B,\Sigma|_{B} \right)$ is standard Borel.
\end{theorem}
\begin{theorem}[Riesz Representation Theorem]
  Let $K$ be a compact Hausdorff space, and $L\colon C(K)\rightarrow \C$ a continuous linear functional. Then, there is a unique complex-valued Borel measure such that
  \begin{align*}
    L(f) &= \int_{}^{} f\:d\mu.
  \end{align*}
  Furthermore, this measure is a Radon measure.
\end{theorem}
\begin{remark}
  From the Riesz Representation Theorem, it can be shown that the following hold:
  \begin{itemize}
    \item $L$ is positive if and only if $\mu$ is a positive measure;
    \item $L(1) = 1$ if and only if $\mu$ is a probability measure.
  \end{itemize}
\end{remark}
One of the fundamental constructions in ergodic theory, especially when we are concerned with actions by a countable discrete group on a standard probability space, is the projective limit of a (countable) family of standard probability spaces. This is a vital theorem.

Consider a system of standard probability spaces $\left( X_k,\Sigma_k \right)_{k=0}^{\infty}$ with connecting maps $\varphi_{k+1}\colon X_{k+1}\rightarrow X_k$ satisfying $\left( \varphi_{k+1} \right)_{\ast}\mu_{k+1} = \mu_k$.

Then, the projective limit
\begin{align*}
  X_{\infty} &= \varprojlim X_k\\
             &= \set{\left( x_k \right)_{k}\in \prod_{k=1}^{\infty}X_k : \varphi_k\left( x_{k+1} \right) = x_k}
\end{align*}
admits a family of projection maps $\varphi_{k,\infty}\colon X_{\infty}\rightarrow X_k$, where we have $\varphi_{k,\infty}\left( \left( x_n \right)_n \right) = x_k$ such that $\varphi_{k}\circ \varphi_{k+1,\infty} = \varphi_{k,\infty}$.
\begin{center}
  % https://tikzcd.yichuanshen.de/#N4Igdg9gJgpgziAXAbVABwnAlgFyxMJZABgBpiBdUkANwEMAbAVxiRAA0B9YAHR6zAAzHAE8AviDGl0mXPkIoyARiq1GLNl2ABrANRKJUmdjwEiS0iur1mrRB07bJqmFADm8IqEEAnCAFskMhAcCCQAJmt1OxA+eh80AAssbm1SPgFhcUlpEF8AoOpQpAsQBjoAIxgGAAVZUwUQHyw3RJwQKNs2OLoE5NT9dP4hUUNc-MDEUuLESLLK6rqTeTZm1vbOjXsevpSnMQoxIA
\begin{tikzcd}
X_{\infty} \arrow[rd, "{\varphi_{k,\infty}}"] \arrow[d, "{\varphi_{k+1,\infty}}"'] &     \\
X_{k+1} \arrow[r, "\varphi_k"']                                                    & X_k
\end{tikzcd}
\end{center}
The projective limit is universal in the sense that if we have a family of maps $f_{k+1}\colon Y\rightarrow X_{k+1}$ satisfying $\varphi_k\circ f_{k+1} = f_k$, then there is a unique map $f\colon Y\rightarrow X$ such that $\varphi_{k,\infty}\circ f = f_k$ for all $k$.

We define the $\sigma$-algebra on $X_{\infty}$ to be
\begin{align*}
  \Sigma_{\infty} &= \sigma\left( \bigcup_{k=1}^{\infty}\set{\varphi_{k,\infty}^{-1}\left( E_k \right) : E_k\in \Sigma_k} \right).
\end{align*}
Then, the \textbf{Kolmogorov Extension Theorem} says the following.
\begin{theorem}[Kolmogorov Extension Theorem]
  There is a unique probability measure $\mu_{\infty}$ on $\left( X_{\infty},\Sigma_{\infty} \right)$ satisfying
  \begin{align*}
    \mu_{\infty} \left( \varphi_{k,\infty}^{-1}\left( E \right) \right) &= \mu_k\left( E \right)
  \end{align*}
  for all $E\in \Sigma_k$.
\end{theorem}
We will give two proofs: one uses the Carathéodory criterion and is more explicit, while the other uses the Riesz Representation Theorem and is more slick.
\begin{proof}[Proof via the Riesz Representation Theorem]
  Since all the maps $\varphi_k$ are Borel, there are zero-dimensional Polish topologies on each set $X_k$ such that the $\varphi_k$ are continuous; since the $X_k$ are all probability spaces and the measures $\mu_k$ are regular, it follows that each $X_k$ is compact. In particular, we also have that the maps $\varphi_{k,\infty}$ are continuous.

  Using Gelfand duality, it follows that for each $X_k$, there is an embedding $C\left( X_k \right)\rightarrow C\left( X_{\infty} \right)$ such that for any $f\in C\left( X_k \right)$, we obtain a continuous map $f\circ \varphi_{k,\infty}\colon X_{\infty}\rightarrow \C$. The set
  \begin{align*}
    A &= \bigcup_{k\geq 1} C\left( X_k \right),
  \end{align*}
  viewing each $C\left( X_k \right)$ as a subset of $C\left( X_{\infty} \right)$, is a unital $\ast$-subalgebra that separates the points of $X_{\infty}$. We may define the linear functional $L\colon A\rightarrow \C$ given by taking
  \begin{align*}
    L|_{C\left( X_k \right)} &= \int_{}^{} \cdot\:d\mu_k.
  \end{align*}
  Notice that this is a positive linear functional satisfying $L(1) = 1$, so since $L$ is continuous on the unital $\ast$-subalgebra $A$, it follows that there is a unique extension to the closure $ L_{\infty}\colon \overline{A} = C\left( X_{\infty} \right)\rightarrow \C $ satisfying $L_{\infty}(1) = 1$. By the Riesz Representation Theorem, there is a unique measure $\mu_{\infty}$ such that
  \begin{align*}
    L_{\infty}\left( f \right) &= \int_{}^{} f\:d\mu_{\infty}
  \end{align*}
  and satisfies
  \begin{align*}
    L_{\infty}|_{C\left( X_k \right)} &= \int_{}^{} \cdot\:d\mu_k.
  \end{align*}
\end{proof}
\begin{proof}[Proof via Carathéodory Criterion]
  Without loss of generality, we may assume that the $X_k$ are compact and the $\varphi_{k+1}$ are Borel. Consider the algebra
  \begin{align*}
    \mathcal{A} &= \bigcup_{k\geq 1} \set{\varphi_{k,\infty}^{-1}\left( E \right) : E\in \Sigma_k}.
  \end{align*}
  Notice that this is the algebra of subsets that generates $\Sigma_{\infty}$. Furthermore, there is a unique set-map $\mu_{\infty}\colon \mathcal{A}\rightarrow [0,1]$ taking $\mu_{\infty}\left( \varphi_{k,\infty}^{-1}\left( E \right) \right) = \mu_k\left( E \right)$ for all $E\in \Sigma_k$. Thus, it suffices to show that $\mu_{\infty}$ is a premeasure.

  It suffices to show that for all decreasing sequences $\left( F_n \right)_n$ with $\mu_{\infty}\left( F_n \right)\geq \ve > 0$, then the intersection of all the $F_n$ is nonempty. This follows from the fact that if one has an ascending sequence $E_n\nearrow E$, then $\left( E\setminus E_n \right)\searrow \emptyset$, so if one takes the Carathéodory criterion and takes the contrapositive of this reversal, we obtain the desired result.

  Let
  \begin{align*}
    F_k &= \varphi_{n_k,\infty}^{-1} \left( E_{n_k} \right).
  \end{align*}
  If $n_k$ does not diverge to $\infty$, then there exists some subsequence $n_{k_j}$ that is constant equal to $n$, and $F_{k_j}$ all have $\varphi_{n,\infty}\left( F_{k_j} \right) = E_{n_{k_j}}$. Since $\mu_n$ is a measure, we are done in this case, following from the fact that all the $F_n$ past $F_{k_j}$ lie in $X_n$.
\end{proof}
The most important case is when the projective limit is the infinite product measure. If $\left( X_n,\mu_n \right)$ is a countable family of standard probability spaces, then the product measure $\bigotimes_{n=1}^{\infty}\mu_n$ is a probability measure.
\subsection{Canonical Examples}%
\begin{example}[Bernoulli Actions]
  Let $\left( X,\mu \right)$ be a standard probability space, and let $G$ be any countable group. Then, there is a unique probability measure on the product space $\left( X,\mu \right)^{G}$, and we have that $G\curvearrowright \left( X,\mu \right)^{G}$ in a probability measure-preserving manner by translation --- that is, $h\cdot \left( x_g \right)_{g\in G} = \left( x_{h^{-1}g} \right)_{g\in G} $. This follows from the fact that $G$ is probability measure-preserving on all the Borel cylinder sets
  \begin{align*}
    \bigcup_{\substack{F\subseteq G\\F\text{ finite}}} \set{\pi_{F}^{-1}\left( E \right) : E\subseteq X^{F}},
  \end{align*}
  so we may apply the monotone class lemma (that no one remembers from measure theory) to see that it applies for all the Borel sets of $\left( X,\mu \right)^{G}$.
\end{example}
\begin{example}
  Suppose $G$ is a compact Hausdorff topological group. Then, there is a unique (regular) Borel probability measure $\mu$, known as the Haar (probability) measure that is invariant under left translation. If $\pi\colon \Gamma\rightarrow G$ is a group homomorphism of a discrete group $\Gamma$ on $G$, then the action $\gamma\cdot g = \pi(\gamma)g$ yields a probability measure-preserving action of $\Gamma$ on $G$.

  For instance, this yields actions $\Z\curvearrowright S^1$ given by rotations --- that is, $e^{2\pi i \theta}\mapsto e^{2\pi i \left( \theta + \alpha \right)}$.

  In general, if $G$ is a compact group, and $\phi\colon \Gamma\rightarrow \aut(G)$ is an action of a discrete group $\Gamma$ that preserves the Haar measure, we call such an action an \textit{algebraic action}.
\end{example}
\begin{example}
  If $M$ is a finite volume Riemannian manifold with volume form $\operatorname{vol}$, and $\phi\in \operatorname{Diff}\left( M \right)$ is a diffeomorphism satisfying $\det\left( \phi'(p)^{\ast}\phi'(p) \right) = 1$ for all $p\in M$, then the pushforward measure satisfies $\phi_{\ast}\operatorname{vol} = \operatorname{vol}$.

  Group actions $\Gamma\to \operatorname{Diff}\left( M \right)$ appear in geodesic and \href{https://en.wikipedia.org/wiki/Horocycle}{horocycle flow}.
\end{example}
\section{The Ergodic Theorems}%
We will now discuss the mean and pointwise ergodic theorems.
\subsection{Mean Ergodic Theorem}%
\begin{proposition}
  Let $X$ be a Banach space. Let $\left( T_n \right)_n\in B(X)$ be a uniformly operator-norm bounded. Then, for all $S\in B(X)$, the sets
  \begin{align*}
    V &= \set{x\in X : T_nx\text{ converges}}\\
    W &= \set{x\in X : \norm{T_nx - Sx}\rightarrow 0}
  \end{align*}
  are norm-closed subspaces of $X$.
\end{proposition}
\begin{theorem}[Mean Ergodic Theorem]
  Let $H$ be a Hilbert space, and let $T\in B(H)$ be such that $\norm{T}\leq 1$. Define
  \begin{align*}
    V &= \set{\xi\in H : T\xi = \xi},
  \end{align*}
  and let $P = \proj_V$. Then, 
  \begin{align*}
    \left( \frac{1}{n}\sum_{j=0}^{n-1}T^{j} \right) &\xrightarrow{\operatorname{SOT}} P.
  \end{align*}
\end{theorem}
\begin{proof}
  Let $K_1 = \img\left( 1-T \right)$ and $K_2 = \ker\left( 1-T \right)$. We claim that $K_1$ and $K_2$ are orthogonal to each other and satisfy $ \overline{K_1} + K_2 = H $. For the first point, we have for all $v\in K_1$ and $w\in K_2$, we write $v = \left( 1-T \right)\xi$, and have
  \begin{align*}
    \iprod{\left( 1-T \right)\xi}{w} &= \iprod{\xi}{w} - \iprod{\xi}{T^{\ast}w}\\
    \norm{T^{\ast}w - w}^2 &= \norm{T^{\ast}w}^2 - 2\re\left( \iprod{w}{Tw} \right) + \norm{w}^2\\
                           &= \norm{T^{\ast}w}^2 - \norm{w}^2\\
                           &\leq 0,
  \end{align*}
  where we note that $\norm{T^{\ast}} = \norm{T}$. Therefore, $T^{\ast}w = w$, so we have that $K_2\subseteq K_1^{\perp}$.

  Now, suppose $\xi\in K_1^{\perp}$. Then,
  \begin{align*}
    \norm{\left( 1-T \right)\xi}^2 &= \norm{T\xi}^2 - 2\re \left( \iprod{T\xi}{\xi} \right) + \norm{\xi}^2,
      \intertext{with}
    \iprod{T\xi}{\xi} &= \iprod{\left( T-1 \right)\xi}{\xi} + \norm{\xi}^2\\
                      &= \norm{\xi}^2
  \end{align*}
  since $\xi\in K_1^{\perp}$. Therefore,
  \begin{align*}
    \norm{\left( 1-T \right)\xi}^2 &\leq \norm{\xi}^2 - \norm{\xi}^2\\
                                   &= 0
  \end{align*}
  since $\norm{T}\leq 1$. Therefore, $K_2 = K_1^{\perp}$.

  Therefore, it suffices to show that
  \begin{align*}
    \frac{1}{n}\sum_{j=0}^{n-1}T^{j}\xi &\rightarrow P\xi
  \end{align*}
  for all $\xi\in K_1\cup K_2$. For each $\xi\in K_2$, we obtain the desired result since both sides are equal to $\xi$.

  For each $\xi\in K_1$, we have that $P\xi = 0$. If $\xi = \left( 1-T \right)\eta$, then
  \begin{align*}
    \frac{1}{n}\sum_{j=0}^{n-1}T^{j}\xi &= \frac{1}{n}\left( 1-T^{n} \right)\eta,
  \end{align*}
  so that
  \begin{align*}
    \norm{\frac{1}{n}\sum_{j=0}^{n-1}T^{j}\xi} &\leq \frac{2}{n}\norm{\eta}\\
                                             &\xrightarrow{n\to\infty}0.
  \end{align*}
\end{proof}
\subsection{Pointwise Ergodic Theorem}%
Now, consider a probability space $\left( X,\mu \right)$ and a measure-preserving map $T\colon X\rightarrow X$. Then, the operator $V_T\in B\left( L^2\left( X,\mu \right) \right)$ is given by $V_Tf = f\circ T$ is an isometry.

If $f\in L^{\infty}\left(X,\mu\right)$, and $P = \proj\left( \operatorname{fix}\left( V_T \right) \right)$, then we have $\norm{Pf}_{L^{\infty}}\leq \norm{f}_{L^{\infty}}$. This follows from the fact that we may find a subsequence such that
\begin{align*}
  \frac{1}{n_k}\sum_{j=0}^{n_k-1}f\circ T^{j} &\rightarrow Pf
\end{align*}
pointwise almost everywhere. Therefore, we have
\begin{align*}
  \norm{Pf}_{L^{\infty}} &\leq \liminf_{k\rightarrow\infty} \frac{1}{n_k}\norm{\sum_{j=0}^{n_k-1}f\circ T^{j}}_{L^{\infty}}\\
                         &\leq \norm{f}_{L^{\infty}}.
\end{align*}
A similar result holds on $L^p$ for general $p$. We discuss this below.
\begin{corollary}
  Suppose $\left( X,\mu \right)$ is a probability space, and $T\colon X\rightarrow X$ is a probability measure-preserving action. Then, for all $f\in L^p\left( X,\mu \right)$ with $1\leq p < \infty$, there exists $ \widetilde{f}\in L^p\left( X,\mu \right) $ such that
  \begin{align*}
    \norm{\widetilde{f} - \frac{1}{n}\sum_{j=0}^{n-1}f\circ T^{j}}_{L^p} &\xrightarrow{n\to \infty}0.
  \end{align*}
  The element $ \widetilde{f} $ satisfies:
  \begin{enumerate}[(i)]
    \item $\widetilde{f}\circ T = \widetilde{f}$ almost everywhere;
    \item $ \ds \int_{}^{} \widetilde{f}g\:d\mu = \int_{}^{} fg\:d\mu $ for all $g\in L^{q}\left( X,\mu \right)$ satisfying $g\circ T = g$ almost everywhere, where $q$ is the conjugate exponent for $p$;
    \item $\norm{\widetilde{f}}_{L^p}\leq \norm{f}_{L^p}$.
  \end{enumerate}
\end{corollary}
\begin{proof}
  In the following proof, we will still be concerned with the operator $P$ defined as
  \begin{align*}
    P &= \proj\left( \operatorname{fix}\left( V_T \right) \right)\\
      &= \proj\left( \set{f\in L^2\left( X,\mu \right) : V_Tf = f} \right).
  \end{align*}
  \begin{enumerate}[(i)]
    \item Let $V_T\in B\left( L^p\left( X,\mu \right) \right)$ be defined by $V_Tf = f\circ T$. Then,
      \begin{align*}
        \norm{\frac{1}{n}\sum_{j=0}^{n-1}V_T^{j}} &\leq 1.
      \end{align*}
      We have that
      \begin{align*}
        W &= \set{f\in L^p\left( X,\mu \right) : \frac{1}{n}\sum_{j=0}^{n-1}V_T^{j}f\text{ converges}}
      \end{align*}
      is a norm-closed subspace. It suffices to show that $L^{\infty}\left( X,\mu \right)\subseteq W$, as then the $L^p$-closure of $L^{\infty}\left( X,\mu \right)$ is contained in $W$, but that is equal to $L^{p}\left( X,\mu \right)$.

      Fix $f\in L^{\infty}\left( X,\mu \right)$.

      If $1\leq p \leq 2$, we have that
      \begin{align*}
        \norm{\frac{1}{n}\sum_{j=0}^{n-1}V_T^{j}f - Pf}_{L^p} &\leq \norm{\frac{1}{n}\sum_{j=0}^{n-1}V_T^{j}f - Pf}_{L^2}\\
                                                              &\rightarrow 0.
      \end{align*}
      If $p\geq 2$, then
      \begin{align*}
        \norm{\frac{1}{n}\sum_{j=0}^{n-1}V_T^{j}f - Pf}_{L^p}^{p} &\leq 2\norm{f}_{L^{\infty}}^{p-2} \norm{\frac{1}{n}\sum_{j=0}^{n-1}V_T^{j}f - Pf}_{L^2}^2\\
                                                                  &\rightarrow 0.
      \end{align*}
    \item We have that $\left( fg \right)\circ T^{j} = \left( f\circ T^{j} \right)g$, so that
      \begin{align*}
        \norm{ \widetilde{f}g - \frac{1}{n}\sum_{j=0}^{n-1}\left( f\circ T^{j} \right)g }_{L^1} &\leq \norm{g}_{q} \norm{ \widetilde{f} - \frac{1}{n}\sum_{j=0}^{n-1}V_T^{j}f }_{L^p}.
      \end{align*}
      In particular, we may write $\widetilde{f}g$ using the asymptotic form for $\widetilde{f}$. Therefore, we have
      \begin{align*}
        \int_{}^{} \widetilde{f}g\:d\mu &= \lim_{n\rightarrow\infty} \frac{1}{n}\sum_{j=0}^{n-1} \int_{}^{} f\left( T^{j}(x) \right)g(x)\:d\mu(x)\\
                                        &= \lim_{n\rightarrow\infty} \frac{1}{n}\sum_{j=0}^{n-1} \int_{}^{} \left( fg \right)\left( T^{j}(x) \right)\:d\mu(x)\\
                                        &= \int_{}^{} fg\:d\mu.
      \end{align*}
    \item We have that
      \begin{align*}
        \norm{\widetilde{f}}_{L^p} &= \lim_{n\rightarrow\infty} \frac{1}{n}\norm{\sum_{j=0}^{n-1}f\circ T^{j}}_{L^p}\\
                                   &\leq \lim_{n\rightarrow\infty}\frac{1}{n}\sum_{j=0}^{n-1} \widetilde{f}\circ T^{j}_{L^p}\\
                                   &= \norm{f}_{L^p}.
      \end{align*}
  \end{enumerate}
\end{proof}
The next result we will be concerned with is the Pointwise Ergodic Theorem. To prove it, we must prove a version of this theorem known as the Maximal Ergodic Theorem. This concerns a maximal function defined by
\begin{align*}
  f^{\ast}\left( x \right) &= \sup_{n} \frac{1}{n}\sum_{j=0}^{n-1} f\circ T^{j}(x).
\end{align*}
\begin{theorem}
  If $f\in L^1\left( X,\mu \right)$, then
  \begin{align*}
    \int_{\set{f^{\ast} > 0}}^{} f\:d\mu &\geq 0.
  \end{align*}
\end{theorem}
\begin{proof}
  The set on which $f^{\ast}$ is strictly greater than zero is the disjoint union of the sets
  \begin{align*}
    B_1 &= \set{x : f(x) \geq 0}\\
        &\vdots\\
    B_n &= \set{x : f(x)\leq 0,f(x) + f\circ T(x) \leq 0,\dots,f(x) + \cdots + f\circ T^{n-2}(x) \leq 0, f(x) + \cdots + f\circ T^{n-1}(x) > 0}.
  \end{align*}
  If we are able to show that
  \begin{align*}
    \int_{B_1\cup\cdots\cup B_n}^{} f\:d\mu &\geq 0
  \end{align*}
  holds, then dominated convergence yields the desired result.

  In order to obtain this, we will break up the set $B_1\cup\cdots\cup B_n$ into a family of disjoint pieces on which the integral of $f$ is strictly positive. We will view them as towers of the form $B_k'\cup T\left( B_k' \right)\cup\cdots\cup T^{k-1}\left( B_k' \right)$.

  First, we observe that $T^{k}\left( B_n \right)\subseteq B_1\cup\cdots\cup B_{n-k}$ for each $k=1,\dots,n-1$, since if $x\in B_n$, then we have
  \begin{align*}
    0\geq f\left( x \right) + f\circ T(x) + \cdots + f\circ T^{k-1}(x)\\
      \intertext{while}
      0 &< f(x) + f\circ T(x) + \cdots + f\circ T^{k-1}(x) + f\circ T^{k}(x) + \cdots + f\circ T^{n-1}(x).
  \end{align*}
  In particular, this means that
  \begin{align*}
    0 &< f\circ T^{k}(x) + \cdots + f\circ T^{n-1}(x)\\
      &= f\left( T^{k}(x) \right) + f\circ T\left( T^{k}(x) \right) + \cdots + f\circ T^{n-k-1}\left( T^{k}(x) \right).
  \end{align*}

  Next, we note that the sets $B_n,T\left( B_n \right),\dots,T^{n-1}\left( B_n \right)$ are pairwise disjoint. This follows from the fact that if $T^{i}\left( B_n \right)\cap T^{j}\left( B_n \right)\neq \emptyset$ for some $i < j$, then $B_n\cap T^{j-i}\left( B_n \right)\neq \emptyset$, but we have already established the family of sets $B_i$ are disjoint, while $T^{j-i}\left( B_n \right)\subseteq B_1\cup\cdots\cup B_{n-(j-i)}$.

  We will define the decomposition of $B_1\cup\cdots\cup B_{n}$ as follows. Take
  \begin{align*}
    B_n' &= B_n\\
    C_n &= B_n'\cup T\left( B_n' \right)\cup\cdots\cup T^{n-1}\left( B_n' \right)\\
    B_{n-1}' &= B_{n-1}\setminus C_n\\
    C_{n-1} &= B_{n-1}'\cup T\left( B_{n-1}' \right) \cup\cdots\cup T^{n-2}\left( B_{n-1}' \right)\\
            &\vdots\\
    B_1' &= B_1\setminus \left( C_2\cup\cdots\cup C_n \right)\\
    C_1 &= B_1'.
  \end{align*}
  The picture one must have in mind is viewing each $C_k$ as a ``column'' in a table with disjoint ``levels'' given by $T^{j}\left( B_k' \right)$ for each $0\leq j \leq k-1$. Each ``base'' $B_k'$ is disjoint 
\end{proof}

\end{document}
